\chapter{Listado de símbolos y convenciones}\label{listado-de-simbolos-y-convenciones}

Este apéndice reúne las siglas, los símbolos matemáticos y las convenciones de
notación empleadas a lo largo del documento. Su finalidad es que cualquier
lector pueda interpretar sin ambigüedad las tablas, ecuaciones y métricas
presentadas en los capítulos anteriores.

\section*{Siglas y acrónimos}
\addcontentsline{toc}{section}{Siglas y acrónimos}

\begin{longtable}[]{@{}ll@{}}
\toprule
\textbf{Sigla} & \textbf{Significado} \\
\midrule
\endfirsthead
\toprule
\textbf{Sigla} & \textbf{Significado} \\
\midrule
\endhead
\bottomrule
\endlastfoot
API & \eng{Application Programming Interface} \\
ARIMA & \eng{AutoRegressive Integrated Moving Average} \\
AUC & \eng{Area Under the Curve} \\
CatBoost & \eng{Categorical Boosting} \\
CCAD & Centro de Computación de Alto Desempeño (UNC) \\
CPU & \eng{Central Processing Unit} \\
DeepAR & \eng{Deep AutoRegressive} \\
EFB & \eng{Exclusive Feature Bundling} \\
ERIA & Ética, Regulaciones e Impacto de la IA (área de LIDeSIA) \\
ETS & \eng{Error, Trend, Seasonality} \\
FCEFyN & Facultad de Ciencias Exactas, Físicas y Naturales \\
GBM & \eng{Gradient Boosting Machine} \\
GNN & \eng{Graph Neural Network} \\
GOSS & \eng{Gradient-based One-Side Sampling} \\
GPU & \eng{Graphics Processing Unit} \\
IA & Inteligencia artificial \\
IEC & \eng{International Electrotechnical Commission} \\
ISO & \eng{International Organization for Standardization} \\
ITID & Infraestructura Tecnológica e Ingeniería de Datos (área de LIDeSIA) \\
LIDeSIA & Laboratorio de Investigación y Desarrollo de Software e IA \\
LightGBM & \eng{Light Gradient Boosting Machine} \\
LLM & \eng{Large Language Model} \\
LSTM & \eng{Long Short-Term Memory} \\
MAE & \eng{Mean Absolute Error} \\
MASE & \eng{Mean Absolute Scaled Error} \\
MIA & Matemática para la IA (área de LIDeSIA) \\
ML & \eng{Machine Learning} (aprendizaje automático) \\
MLP & \eng{Multilayer Perceptron} (perceptrón multicapa) \\
MPF & Ministerio Público Fiscal de la Provincia de Córdoba \\
MSE & \eng{Mean Squared Error} \\
MVP & \eng{Minimum Viable Product} (producto mínimo viable) \\
PatchTST & \eng{Patch Time Series Transformer} \\
PMI & \eng{Project Management Institute} \\
PPT & Personas, procesos y tecnología \\
PR-AUC & \eng{Precision-Recall Area Under the Curve} \\
RAG & \eng{Red, Amber, Green} (codificación de alertas por color) \\
RAM & \eng{Random Access Memory} \\
RMSE & \eng{Root Mean Squared Error} \\
ROC & \eng{Receiver Operating Characteristic} \\
SAC & Sistema de Administración de Causas \\
SAT & Sistema de alerta temprana \\
SGAV & Sistema de Gestión de Siniestros Viales \\
SLURM & \eng{Simple Linux Utility for Resource Management} \\
TFT & \eng{Temporal Fusion Transformer} \\
TRL & \eng{Technology Readiness Level} (nivel de madurez tecnológica) \\
UJ & Unidad Judicial \\
UNC & Universidad Nacional de Córdoba \\
XAI & \eng{Explainable Artificial Intelligence} (IA explicable) \\
XGBoost & \eng{Extreme Gradient Boosting} \\
\end{longtable}

\section*{Símbolos matemáticos}
\addcontentsline{toc}{section}{Símbolos matemáticos}

\begin{longtable}[]{@{}>{$}l<{$}p{10.5cm}@{}}
\toprule
\multicolumn{1}{@{}l}{\textbf{Símbolo}} & \textbf{Significado} \\
\midrule
\endfirsthead
\toprule
\multicolumn{1}{@{}l}{\textbf{Símbolo}} & \textbf{Significado} \\
\midrule
\endhead
\bottomrule
\endlastfoot
K & Cantidad de variables medidas simultáneamente en el sistema multivariado. \\
T & Cantidad total de instantes de tiempo de la serie. \\
t & Índice del paso de tiempo (bloque de seis horas). \\
\mathbf{z}_t & Vector de observaciones en el instante $t$, con $\mathbf{z}_t \in \mathbb{R}^{K}$. \\
\mathbf{Z} & Matriz de datos original, de dimensiones $T \times K$. \\
n & Cantidad de observaciones del conjunto de entrenamiento. \\
N & Cantidad de muestras del conjunto de prueba. \\
y_i & Valor observado de la variable objetivo para la observación $i$. \\
\hat{y}_i & Valor predicho por el modelo para la observación $i$. \\
\hat{y}^{(t-1)}_i & Predicción acumulada tras $t-1$ iteraciones de \eng{boosting}. \\
f_t & Árbol incorporado en la iteración $t$ del \eng{boosting}. \\
\mathcal{L}^{(t)} & Función objetivo de XGBoost en la iteración $t$. \\
\Omega(f_t) & Término de regularización que penaliza la complejidad del árbol. \\
g_i & Gradiente de primer orden de la función de pérdida. \\
h_i & Gradiente de segundo orden (hessiano) de la función de pérdida. \\
M & Cantidad de modelos que participan del ensamble ponderado. \\
w_m & Peso asignado al modelo $m$ en el ensamble, con $\sum w_m = 1$ y $w_m \geq 0$. \\
m & Período estacional de la serie ($m = 4$ bloques por día). \\
z & Puntaje estándar (\eng{Z-Score}) usado por la capa de decisión. \\
\sigma & Desviación estándar del comportamiento histórico. \\
\text{P10},\ \text{P90} & Cuantiles 10\,\% y 90\,\% del intervalo de predicción. \\
\end{longtable}

\section*{Convenciones de notación}
\addcontentsline{toc}{section}{Convenciones de notación}

\begin{description}
  \item[Separadores numéricos.] Se emplea la coma como separador decimal y el
    punto como separador de miles a partir de cinco cifras, según la convención
    del español rioplatense. Los valores tomados de salidas de software fueron
    convertidos a esta notación.
  \item[Nombres de variables.] Los campos del conjunto de datos y los
    hiperparámetros de los modelos se escriben en tipografía monoespaciada,
    respetando la grafía exacta con la que figuran en el código:
    \var{TARGET\_CRITICO}, \var{promedio\_victimas\_24h}, \var{max\_depth}.
  \item[Términos en inglés.] Los anglicismos técnicos sin traducción consolidada
    se escriben en cursiva (\eng{forecasting}, \eng{baseline}, \eng{recall}).
    Cuando existe equivalente en español se usa este último, indicando el
    término original entre paréntesis en su primera aparición.
  \item[Nombres de modelos y librerías.] Se respeta la grafía oficial de cada
    proyecto: XGBoost, LightGBM, AutoGluon, AutoGluon-TimeSeries, CatBoost.
  \item[Tablas y figuras.] Los elementos tabulares se denominan \emph{tablas} y
    los gráficos, \emph{figuras}; ambos se numeran por capítulo y se referencian
    de forma cruzada. Cada uno lleva al pie una nota en formato APA que indica
    su alcance y su fuente.
  \item[Niveles de alerta.] La codificación por color sigue el criterio definido
    en la sección de tipos de alerta: verde (línea de base), amarillo
    (prevención), naranja (alerta) y rojo (crítico).
  \item[Bloques temporales.] Salvo indicación en contrario, la unidad de
    agregación es el bloque de seis horas; el bloque 1 abarca de 00:00 a 06:00 y
    los restantes se numeran de forma consecutiva.
\end{description}
